% Pouzil som prostredie python 
% kniznicu BioPython, ktora poskytuje rozne funkcie A na lokalne vyhlavadave som pouzil Bio.align.
% Tiez som Pouzil tkinter na filedialog a messagebox, aby som mohol jednoducho vyberat subory a zobrazovat spravy.
% A kniznicu PySam pre robenie s Bam subormi.

%Todo: zformuj do akademickeho neosobneho stylu vsetko co je v hore
\subsection{Python a BioPython}
Python je vysokoúrovňový programovací jazyk, ktorý sa vyznačuje jednoduchou syntaxou a širokou škálou knižníc, čo ho činí ideálnym nástrojom pro bioinformatické analýzy.

BioPython je rozsiahla knižnica pre biologické výpočty v Pythone, ktorá poskytuje nástroje na prácu s biologickými dátami, vrátane sekvencií DNA, RNA.
V tejto práci bol Python použitý na implementáciu rôznych krokov spracovania a analýzy DNA sekvencií, vrátane lokalizácie primerov, klastrovania a viacnásobného zarovnania.
BioPython bol použitý hlavne pre vyhľadávanie primerov a zarovnanie sekvencií pomocou modulu Bio.Align.
Knižnica PySam bola použitá pre prácu s BAM súbormi, ktoré vznikajú po preklade nanopórových čítaní nástrojom Dorado.
Z knižnice tkinter bolo použité filedialog pre vytvorenie jednoduchého grafického rozhrania na výber súborov, čo zjednodušilo interakciu počas spracovania dát.

